\section{Metadata-only large molecular datasets with native partitions}

PCQM4Mv2, MOSES and ogbg-molpcba extend the application evidence beyond ODAC25
and the nested ChEMBL scaling samples. The metered archive/CSV payload is
\NativeMetadataMB{} MB, including ZIP catalog ranges; small discovery pages,
source JSON and HTTP headers are outside this tally. Retrieval is capped at
200 MiB of dataset payload and 128 MiB per object. Only the required SMILES,
mapping and native-index ZIP members are fetched. No coordinates, SDFs,
trajectory archives, processed graph arrays or model checkpoints are downloaded.

MOSES files are pinned to an immutable maintainer commit and checked against
their Git-LFS SHA-256 identities. Archive-member CRCs, extracted-file hashes,
range requests and response sizes are retained for OGB sources. Raw metadata
and locally generated packed fingerprints are reproducible caches outside the
public evidence bundle. The CSV property and assay columns are not loaded for
optimization. Sources and recipes are in \path{revision_20261104_native_splits}.

\subsection{Native roles, identities and parsing}

The original source assignments are checked for disjointness and complete row
coverage before fingerprinting. PCQM4Mv2 retains its train, valid, test-dev and
test-challenge roles; the released split is CID-based, as specified in the
\href{https://ogb.stanford.edu/docs/lsc/pcqm4mv2/}{OGB documentation}, rather
than the originally intended scaffold split. MOSES retains train, random test
and scaffold test. Ogbg-molpcba retains train, valid and test under its native
scaffold assignment. No official held-out partitions are merged.
Released PCQM4Mv2 test-dev and test-challenge response labels are unavailable.
Its repartitioned assignments support structural and cost analysis; using them
as a labelled prediction benchmark would require the corresponding responses.

The same radius-two, 1,024-bit Morgan fingerprint recipe is used throughout.
RDKit 2026.03.5 cannot parse 39 PCQM4Mv2 records and two ogbg-molpcba records;
their original positions and failure reasons are recorded. Every method uses
the same retained rows and native count vector. All MOSES rows are retained.
The canonical-identity audit uses hashes of isomeric canonical SMILES. Repeated
records are retained, and neither native nor optimized assignments are silently
deduplicated. In PCQM4Mv2, distinct source records can share a canonical graph;
this does not establish identical atomistic geometries or a predictive effect.

\begin{longtable}{lrrrrr}
\caption{Metadata identity audit with the recorded RDKit version. Keys hash isomeric canonical SMILES; no chemical deduplication is silently applied. Crossing keys occur in more than one original native role. Parsing failures are fixed before all methods are compared, with original-row maps retained. Counts are molecular records, not atomistic frames.}\\
\toprule
Dataset & Original & Retained & Failed & Unique keys & Crossing keys\\\midrule
\endfirsthead
\toprule
Dataset & Original & Retained & Failed & Unique keys & Crossing keys\\\midrule
\endhead
\midrule\endfoot
\bottomrule\endlastfoot
PCQM4Mv2 & 3,746,620 & 3,746,581 & 39 & 3,574,491 & 34,992\\
MOSES & 1,936,962 & 1,936,962 & 0 & 1,936,962 & 0\\
ogbg-molpcba & 437,929 & 437,927 & 2 & 437,927 & 0\\
\end{longtable}


The numerical package remains pinned to \texttt{f7e0866}. Its ordinary
\texttt{SplitSpec} wrapper permits two or three roles. A separately hashed
\texttt{NativeSplitSpec} subclass validates the four named PCQM4Mv2 roles and
passes them to the unchanged k-way numerical routines. Exact/corridor small
four-fold cases and invalid specifications are regression-tested. The initial
API preflight failed during warm-up and is preserved separately; it produced
no large optimized partitions and is not counted as a scientific outcome.

\subsection{Generation cost and quality measurements}

The production grid has \NativeCompletedCases{} completed cases and
\NativeBudgetOutcomes{} whole-case budget outcomes: one native, three
count-matched random, three exact-count assignment-only and three corridor
plain cases per dataset. Low-rank settings use 256 landmarks, four starts and
at most 25 assignment iterations. Exact-count controls disable refinement;
the 1\% corridor setting refines two candidates when the unchanged size rule
permits it. Above two million rows, point refinement is skipped. The two
settings vary both feasibility and refinement work and do not isolate the
effect of tolerance alone.

Fresh processes use CPUs 24--39 on lambda2, with CUDA hidden and native pools
capped at 16. The measured public call includes its representation construction,
assignment and enabled refinement. Imports, prepared-input loading, unpacking,
input-adapter construction, common warm-up and independent auditing are outside
that clock. Peak RSS is read before auditing and includes inputs/framework costs.
Fingerprint preparation and metadata loading are recorded separately. Published
native construction times and memory were not supplied, so native assignment
loading is never presented as a competing optimization measurement. A
1,200-second whole-case cap includes every stage, and budget outcomes receive
no imputed split-call time.

All methods share 262,144 uniformly sampled ordered pairs of distinct retained
records per dataset. Original Tanimoto is evaluated exactly from packed bits.
Global overlap is normalized by the analytic random crossing fraction at each
returned count vector. Approximate 95\% Monte Carlo intervals resample 64
independent pair blocks; they describe sampling uncertainty at fixed partitions,
not future-dataset generalization or simultaneous coverage over every setting.

For every held-out role, 128 uniform queries without replacement are searched
against the complete returned training fold. Exact packed-bit Tanimoto uses
only mathematically valid upper-bound pruning; it is not approximate neighbour
search. Per-query bootstrap intervals describe query-sampling uncertainty. The
main-table aggregate weights role means by their actual populations, while the
following table preserves their separate meanings. A similarity of one denotes
fingerprint identity, not necessarily identical chemical graphs or geometries.

\begin{longtable}{lrrrrrr}
\caption{Per-held-out-role exact nearest-training exposure on 128 uniform sampled queries per role and seed. Every query is compared with every returned training record. Fractions at one measure fingerprint identity, not identical geometries. Rows average completed seeds; full query-bootstrap intervals and per-seed observations accompany the evidence.}\\
\toprule
Dataset & Assignment & Role & Runs & Nearest & $>0.6$ & $=1$\\\midrule
\endfirsthead
\toprule
Dataset & Assignment & Role & Runs & Nearest & $>0.6$ & $=1$\\\midrule
\endhead
\midrule\endfoot
\bottomrule\endlastfoot
PCQM4Mv2 & Native & valid & 1 & 0.704 & 0.711 & 0.133\\
PCQM4Mv2 & Native & test-dev & 1 & 0.703 & 0.727 & 0.133\\
PCQM4Mv2 & Native & test-challenge & 1 & 0.722 & 0.797 & 0.141\\
PCQM4Mv2 & Random & valid & 3 & 0.760 & 0.878 & 0.146\\
PCQM4Mv2 & Random & test-dev & 3 & 0.771 & 0.904 & 0.154\\
PCQM4Mv2 & Random & test-challenge & 3 & 0.764 & 0.888 & 0.159\\
PCQM4Mv2 & Exact assignment & valid & 3 & 0.488 & 0.234 & 0.000\\
PCQM4Mv2 & Exact assignment & test-dev & 3 & 0.548 & 0.341 & 0.000\\
PCQM4Mv2 & Exact assignment & test-challenge & 3 & 0.552 & 0.375 & 0.000\\
PCQM4Mv2 & 1\% corridor & valid & 3 & 0.493 & 0.234 & 0.000\\
PCQM4Mv2 & 1\% corridor & test-dev & 3 & 0.554 & 0.328 & 0.000\\
PCQM4Mv2 & 1\% corridor & test-challenge & 3 & 0.540 & 0.320 & 0.000\\
MOSES & Native & test & 1 & 0.739 & 0.930 & 0.008\\
MOSES & Native & test\_scaffolds & 1 & 0.707 & 0.883 & 0.000\\
MOSES & Random & test & 3 & 0.759 & 0.938 & 0.016\\
MOSES & Random & test\_scaffolds & 3 & 0.755 & 0.945 & 0.010\\
MOSES & Exact assignment & test & 3 & 0.608 & 0.487 & 0.000\\
MOSES & Exact assignment & test\_scaffolds & 3 & 0.601 & 0.510 & 0.000\\
MOSES & 1\% corridor & test & 3 & 0.608 & 0.495 & 0.000\\
MOSES & 1\% corridor & test\_scaffolds & 3 & 0.587 & 0.424 & 0.000\\
ogbg-molpcba & Native & valid & 1 & 0.652 & 0.672 & 0.023\\
ogbg-molpcba & Native & test & 1 & 0.651 & 0.688 & 0.016\\
ogbg-molpcba & Random & valid & 3 & 0.754 & 0.924 & 0.039\\
ogbg-molpcba & Random & test & 3 & 0.761 & 0.922 & 0.039\\
ogbg-molpcba & Exact assignment & valid & 3 & 0.563 & 0.427 & 0.000\\
ogbg-molpcba & Exact assignment & test & 3 & 0.576 & 0.427 & 0.000\\
ogbg-molpcba & 1\% corridor & valid & 3 & 0.559 & 0.409 & 0.000\\
ogbg-molpcba & 1\% corridor & test & 3 & 0.568 & 0.414 & 0.000\\
\end{longtable}

\begin{longtable}{lrrrr}
\caption{Complete native-split outcome inventory. Budget denotes the 1,200-second whole-case limit, including loading, warm-up, optimization and auditing; it is not an imputed split-call runtime. Native fixed-assignment construction times are unknown. Exact assignment explicitly disables refinement; the corridor setting uses the unchanged size-dependent refinement rule.}\\
\toprule
Dataset & Assignment & Seed & Outcome & Call (s)\\\midrule
\endfirsthead
\toprule
Dataset & Assignment & Seed & Outcome & Call (s)\\\midrule
\endhead
\midrule\endfoot
\bottomrule\endlastfoot
PCQM4Mv2 & Native & 0 & Complete & --\\
PCQM4Mv2 & Random & 0 & Complete & 0.1\\
PCQM4Mv2 & Exact assignment & 0 & Complete & 317.4\\
PCQM4Mv2 & 1\% corridor & 0 & Complete & 299.1\\
PCQM4Mv2 & Random & 1 & Complete & 0.1\\
PCQM4Mv2 & Exact assignment & 1 & Complete & 285.0\\
PCQM4Mv2 & 1\% corridor & 1 & Complete & 283.4\\
PCQM4Mv2 & Random & 2 & Complete & 0.1\\
PCQM4Mv2 & Exact assignment & 2 & Complete & 309.7\\
PCQM4Mv2 & 1\% corridor & 2 & Complete & 318.9\\
MOSES & Native & 0 & Complete & --\\
MOSES & Random & 0 & Complete & 0.1\\
MOSES & Exact assignment & 0 & Complete & 185.3\\
MOSES & 1\% corridor & 0 & Complete & 183.2\\
MOSES & Random & 1 & Complete & 0.1\\
MOSES & Exact assignment & 1 & Complete & 194.0\\
MOSES & 1\% corridor & 1 & Complete & 195.7\\
MOSES & Random & 2 & Complete & 0.1\\
MOSES & Exact assignment & 2 & Complete & 178.5\\
MOSES & 1\% corridor & 2 & Complete & 177.1\\
ogbg-molpcba & Native & 0 & Complete & --\\
ogbg-molpcba & Random & 0 & Complete & 0.0\\
ogbg-molpcba & Exact assignment & 0 & Complete & 38.5\\
ogbg-molpcba & 1\% corridor & 0 & Complete & 42.0\\
ogbg-molpcba & Random & 1 & Complete & 0.0\\
ogbg-molpcba & Exact assignment & 1 & Complete & 39.6\\
ogbg-molpcba & 1\% corridor & 1 & Complete & 36.4\\
ogbg-molpcba & Random & 2 & Complete & 0.0\\
ogbg-molpcba & Exact assignment & 2 & Complete & 43.5\\
ogbg-molpcba & 1\% corridor & 2 & Complete & 42.3\\
\end{longtable}


\subsection{Independent verification and returned-cut fidelity}

Five regression tests check packed metrics against RDKit, zero-mass handling,
four-role capacities and specification validation. Independent RDKit checks
1,024 sampled pair values per dataset. A second NumPy implementation verifies
eight saved queries per held-out role against every training entity for native,
exact-count and corridor seed-zero cases: 168 full-training query checks in
total. Every saved assignment, count, query-selection sequence, threshold score,
mean and sampling interval is also replayed from the saved artifacts.

The additional post-hoc fidelity diagnostic reconstructs sampled factor rows
from saved landmark identities using the pinned double-construction and
float32-storage formula. A regression against the actual producer factor
checks the reconstruction. It evaluates both returned assignments for every
seed without rebuilding a full factor or similarity matrix. No rank or method
selection uses these outcomes. The absolute signed cut residual is normalized
by original kernel mass and by the original returned cut. Approximate
block-bootstrap intervals are reported separately from the earlier full-matrix
tests and bounded accuracy-selection policy; they do not certify uniform
fidelity or optimizer optimality.

\begin{longtable}{lrrrr}
\caption{Post-hoc returned-cut fidelity at 256 landmarks. The factor is reconstructed only at uniform sampled pairs using saved landmarks and the pinned double-construction, float32-storage formula. Each cell is the largest estimate or interval endpoint over three seeds. The intervals are approximate Monte Carlo block-bootstrap intervals for the tested returned partition; they certify neither uniform fidelity nor optimizer optimality.}\\
\toprule
Dataset & Setting & Mass error (\%) & Cut error (\%) & Upper CI (\%)\\\midrule
\endfirsthead
\toprule
Dataset & Setting & Mass error (\%) & Cut error (\%) & Upper CI (\%)\\\midrule
\endhead
\midrule\endfoot
\bottomrule\endlastfoot
PCQM4Mv2 & Exact assignment & 0.40 & 3.50 & 3.71\\
PCQM4Mv2 & 1\% corridor & 0.39 & 3.50 & 3.71\\
MOSES & Exact assignment & 0.32 & 1.30 & 1.38\\
MOSES & 1\% corridor & 0.32 & 1.31 & 1.39\\
ogbg-molpcba & Exact assignment & 0.41 & 1.53 & 1.62\\
ogbg-molpcba & 1\% corridor & 0.41 & 1.55 & 1.64\\
\end{longtable}

